%=======================================================================
%  2  CHALLENGE & SCOPING
%  Output of the Day-1 afternoon task (Scoping V6, slides 28 + 29).
%  Required output:  1. context & constraints  2. system view figure with
%  IN/OUT boundary  3. function list / functional diagram  4. KPI table
%  5. final "How might we ...?" question.
%
%  Kept to the five required outputs. The supporting process tables --- the
%  candidate design questions, the full constraint list and the parallel-system
%  comparison --- are in Appendix C, and the decisions they record are in the
%  Decision Log (Appendix B).
%=======================================================================
\section{Challenge \& Scoping}

\subsection{Problem definition and design question}
\label{sec:designquestion}

\textbf{Given challenge (Challenge~1).} \emph{Cities are facing increasing heat
stress due to climate change. Develop bioinspired strategies to reduce urban
heat through building materials, façade design, or urban planning.}

The challenge as given spans three scales --- material, building component and
urban planning. Working on all three would make the project larger rather than
more focused, which the systems view is explicitly meant to prevent. The design
question is the instrument that makes this cut.

\paragraph{Candidates considered.} Five questions were formulated in the team
and assessed against the criterion that a design question must be neither too
broad nor too narrow (Table~\ref{tab:questions}). One was rejected as too broad
(no place, no time, no intervention point), one as unanchored, and one as
factually wrong --- heat waves do not cause heat islands, which are a permanent
consequence of urban surface properties and geometry \parencite{oke1982}. The
two survivors were a night-time removal question, well formed but discarding the
team's central idea of \emph{using} the heat, and a \enquote{turn heat into a
resource} question whose seasonal-storage reading is neither verifiable within
the project nor biologically supported: organisms store chemical energy across
seasons, not sensible heat. The selected question keeps the night-time focus of
the first and the resource framing of the second.

\medskip
\noindent\fbox{\parbox{0.97\textwidth}{%
\textbf{Design question (solution-neutral):}\\[2mm]
\emph{How might we let the heat that Zurich's existing buildings absorb during
a summer day drive its own removal, so that the hotter it gets, the more the
building cools itself?}}}

\medskip
The question names no solution --- no coating, no greening, no geometry, no
water --- while fixing the location, the object (existing buildings, so retrofit
rather than new build) and the season. What it adds to the two surviving
candidates is the subordinate clause, and that clause is the content of the
project: it demands a \textbf{self-regulating} response, in which cooling power
rises with the heat load instead of being fixed. That is what separates the
project from the incumbent measure. A bright surface
(Section~\ref{sec:stateofart}) has a fixed optical response: it reflects the
same fraction at \qty{20}{\celsius} as at \qty{40}{\celsius} and does nothing
once the sun has set. A self-regulating surface behaves like a biological one
--- transpiration in a leaf is driven by the same solar energy that heats the
leaf, so the cooling effort scales with the load, without controller, sensor or
pump. The clause also resolves the team's initial difficulty of carrying two
systems of interest (\enquote{reduce heat} and \enquote{prevent heat}): these
are two functions of one system, prevention acting in the day window and removal
in the night window, coupled by the self-regulating link.

\subsection{Context and constraints}
\label{sec:context}

\paragraph{Reference city and site.} All context values, benchmarks and
engineering checks refer to the City of Zurich (\ang{47.38}\,N, \ang{8.54}\,E,
about \qty{400}{\metre} above sea level); fixing one city is what makes the
constraints and KPI targets below checkable instead of generic. Zurich has an
adopted heat-mitigation strategy, the \emph{Fachplanung Hitzeminderung}
\parencite{stadtzuerich2020}, with an implementation agenda
\parencite{stadtzuerich_agenda}, a leaflet for building owners
\parencite{stadtzuerich_merkblatt} and thematic maps
\parencite{stadtzuerich_teilplaene}, so the project has a real client, a
documented baseline and published target values instead of invented ones; open
geodata \parencite{stadtzuerich_opendata} and MeteoSwiss stations
\parencite{meteoswiss} supply the boundary conditions ($G$, $T_a$, wind,
humidity) for Section~3.2. As reference site we adopt the city's own model area
\emph{Modellierungsgebiet 03 --- Geschlossene Randbebauung}
\parencite[148--151]{stadtzuerich2020}: two outward-closed blocks of five
storeys enclosing courtyards in a dense central district, street space sealed
from façade to façade. The city assesses this type as overheating comparatively
quickly and already publishes modelled baseline and climate-optimised results
for exactly this area --- the methodical advantage --- so our concept can be
compared against a documented reference rather than our own assumptions.
\todo{Record address, orientations, roof area and year of construction of the
building you pick.} The users are the residents of the upper floors, who carry
the strongest load and usually have no mechanical cooling; the pedestrians in
the canyon, who receive façade radiation at eye level; the owner, who pays for the retrofit; and
the trades who install and maintain it.

\paragraph{When --- the design case and its windows.} The design case is not a
single hot afternoon but a \textbf{multi-day heat wave}: five consecutive days
with $T_\mathrm{max} \geq \qty{30}{\celsius}$ and no rain. The concept relies on
a stored resource --- a system that performs on day one and is exhausted on day
three has not solved the problem --- and only a multi-day case puts the water
reserve (F7) and the reset of the heat store (F4) under stress. KPI~6 follows
from it. \todo{Confirm duration and frequency of such spells from MeteoSwiss
records; state station and reference period.} Within each day the design must
perform in the three windows of Table~\ref{tab:windows}, which follow one
another without a gap because the surface never stops exchanging heat. The
remaining hours, 06:00--11:00, are the margin the night window has to leave
behind, so the state of the system at 06:00 is reported as the reset condition
for the next day. The team's initial scoping put the second window at
20:00--24:00; it was moved to 18:00 and split from the night because a window
must contain the instants at which its own KPIs are read (KPI~5 at 02:00), and
because outdoor amenity and sleep quality need different KPIs.

\begin{table}[htbp]
\centering
\caption{Critical operating windows within one day of the design heat wave.
Sunset in Zurich in high summer is around 21:00, sunrise around 06:00.}
\label{tab:windows}
\renewcommand{\arraystretch}{1.2}
\small
\begin{tabularx}{\textwidth}{L{2.2cm} L{3.6cm} X L{1.3cm}}
\toprule
\textbf{Window} & \textbf{What dominates} & \textbf{Why it is critical} & \textbf{KPIs} \\
\midrule
\textbf{Load}\newline 11:00--18:00
& Absorbed short-wave radiation; clear-sky global irradiance
  \qtyrange{900}{950}{\watt\per\square\metre} horizontal around solar noon,
  air up to \qtyrange{33}{36}{\celsius}.
& The energy that has to be dealt with all night enters here; a sealed,
  sun-exposed surface can exceed \qty{60}{\celsius}
  \parencite{stadtzuerich_merkblatt}. The city evaluates PET at 14:00, inside
  this window, so our results stay comparable to Table~\ref{tab:incumbent}.
& 2, 3, 4 \\[1mm]
\textbf{Evening}\newline 18:00--24:00
& Low-angle sun on west-facing surfaces until sunset, then the start of the
  discharge into the canyon.
& The window in which the effect is actually experienced: people are outdoors
  and then go to bed. West façades reach their highest incidence angle here,
  which is why the window starts at 18:00.
& 4, 5 \\[1mm]
\textbf{Night}\newline 00:00--06:00
& Long-wave emission to the sky and continued release of stored heat; clearest
  sky of the cycle.
& Zurich's heat island reaches \qty{7}{\kelvin} at night against
  \qtyrange{1}{2}{\kelvin} by day \parencite{stadtzuerich_merkblatt} --- the
  larger half of the problem, and the half the incumbent does not address. It
  sets a hard deadline: the store must be discharged by sunrise or it saturates
  over the following days.
& 1, 5, 6 \\
\bottomrule
\end{tabularx}
\end{table}

\paragraph{Where on the envelope? A consequence of Zurich's latitude.} Fixing
the city fixes the solar geometry, and that already narrows the design space. At
\ang{47.38}\,N the solar altitude at noon on 21 June is
$h = 90^{\circ} - 47.38^{\circ} + 23.44^{\circ} \approx 66^{\circ}$, so the
direct beam strikes a horizontal roof at an incidence factor of
$\sin h \approx 0.91$ but a vertical south façade at only
$\cos h \approx 0.41$: in midsummer the flat roof receives more than twice the
direct irradiance of a south façade at noon, and it has a far larger sky view
factor for night-time emission. Both facts argue for the roof as a good place to
\emph{act} and against it as the place where the problem is. The heat a roof
stores is released \emph{upward}, to a sky it can see; the heat a façade stores
is released \emph{sideways}, into the canyon, at the height of the people in it.
Since Section~1 fixes the pedestrian-level outdoor space as the place where the
effect has to appear, the system of interest is the \textbf{façade} --- not
because it receives more radiation, but because its heat is released where it
does harm. A façade concept also does not compete with photovoltaics for the
roof area. The roof does not disappear; it changes role, from a second surface
to be treated into the \emph{interface with the good sky view}, which is what
the concept in Section~\ref{sec:creating} requires.
\todo{Confirm the load profile per orientation \parencite{sonnendach}; east and
west façades differ in the shape of the profile, not only its size.}

\paragraph{Constraints.} Table~\ref{tab:constraints} in Appendix~C lists the
context conditions in full and separates hard constraints (a violation
disqualifies a concept) from soft ones (they weight the evaluation in
Section~3.3). Seven are hard: \textbf{no auxiliary electrical energy} for the
cooling function, \textbf{no potable water}, \textbf{retrofit only} onto an
existing envelope, acceptability in the Zurich \textbf{townscape} including
heritage protection, \textbf{no specular glare}, survival of the Swiss
\textbf{annual cycle} to about \qty{-10}{\celsius} with freeze--thaw and UV, and
\textbf{no appreciable increase in heating demand} --- a high solar reflectance
also reduces useful winter gains, and that trade-off must be quantified, not
ignored. Glare and soiling are the two objections the city itself raises against
bright envelopes \parencite[128]{stadtzuerich2020}. The soft constraints cover
installability, maintenance, photovoltaic yield, dead load, embodied greenhouse
gas and cost; the last two were originally KPIs, and Section~\ref{sec:kpi}
explains the move. Installation and maintenance are recorded as constraints
rather than as subsystems, because a system of interest contains components, not
people.

\subsection{State of the art: what the City of Zurich already does}
\label{sec:stateofart}

Scoping is incomplete without knowing what the client already does. Of the
thirteen courses of action (\emph{Handlungsansätze}) in the Fachplanung
Hitzeminderung, four act on the building envelope: HA~09 green roofs, HA~10
green façades, HA~11 high-albedo façade and roof materials, and HA~12 shading of
the outdoor space \parencite[120--131]{stadtzuerich2020}. HA~11 is the direct
incumbent of our design question, and the city already recommends it to building
owners \parencite{stadtzuerich_merkblatt}.

\begin{table}[htbp]
\centering
\caption{Modelled effectiveness of the incumbent measures, from the Fachplanung
Hitzeminderung. Day values are the physiologically equivalent temperature (PET)
at 14:00, \qty{2}{\metre} above ground; night values are air temperature.
\enquote{Range} is the distance over which the effect is measurable.}
\label{tab:incumbent}
\renewcommand{\arraystretch}{1.2}
\small
\begin{tabularx}{\textwidth}{L{4.4cm} c c c X}
\toprule
\textbf{Measure} & \textbf{Median} & \textbf{Max.} & \textbf{Range} & \textbf{Source} \\
\midrule
HA~11 bright façade --- day (PET) & \qty{-1.0}{\kelvin} & \qty{-1.8}{\kelvin} & \qtyrange{2}{3}{\metre} & \multirow{4}{*}{\parencite[Tab.~11, 128]{stadtzuerich2020}} \\
HA~11 bright façade --- night     & \qty{-0.1}{\kelvin} & \qty{-0.2}{\kelvin} & \qtyrange{1}{2}{\metre} & \\
HA~11 bright roof --- day (PET)   & \qty{-1.2}{\kelvin} & \qty{-2.1}{\kelvin} & \qtyrange{2}{4}{\metre} & \\
HA~11 bright roof --- night       & \qty{-0.1}{\kelvin} & \qty{-0.3}{\kelvin} & \qtyrange{1}{2}{\metre} & \\
\midrule
HA~06 high-albedo ground surface instead of asphalt --- day (PET) & \qty{-1.5}{\kelvin} & \qty{-2.8}{\kelvin} & \qtyrange{2}{4}{\metre} & \multirow{2}{*}{\parencite[Tab.~5, 114]{stadtzuerich2020}} \\
HA~06 --- night & \qty{-0.2}{\kelvin} & \qty{-0.5}{\kelvin} & \qtyrange{2}{3}{\metre} & \\
\bottomrule
\end{tabularx}
\end{table}

\paragraph{Baseline values this fixes for us.} The city's model uses a solar
reflectance of $0.3$ for un-greened façades and $0.2$ for asphalt, and
classifies $\leq 0.30$ as low and $\geq 0.80$ as high albedo
\parencite[114, 128--129]{stadtzuerich2020}. Our reference case and the target
of KPI~2 are taken directly from these figures, so the comparison is made on the
client's terms rather than on ours.

\paragraph{The gap --- and the design opportunity.} Table~\ref{tab:incumbent}
contains the central finding for this project. The incumbent works during the
day but is close to ineffective at night: \qty{-0.1}{\kelvin} median air
temperature, against a heat island that reaches \qty{7}{\kelvin}. The reason is
structural, not a question of tuning. A bright surface reduces what is
\emph{absorbed}; it opens no path for releasing heat already stored in the
construction, and its optical response is fixed --- it does not do more when the
load is higher, and nothing at all after sunset. A bio-inspired concept
therefore earns its place if it does at least one of three things: open a
removal path and so deliver a night-time effect; scale its response with the
load, so the worst hours receive the strongest cooling; or resolve glare or
soiling, the two constraints the city names for HA~11 and precisely the points
where biological surfaces are strong.

\paragraph{Two precedents outside the city's own catalogue.} The Fachplanung
lists what the client does, not what the field has built. \emph{StoColor
Lotusan}, a micro-structured superhydrophobic façade paint modelled on the lotus
leaf, has been on the market since 1999 and is stated by its manufacturer to
resist algae and fungal growth \emph{without biocide film protection}
\parencite{sto_lotusan} --- a commercial answer to exactly the soiling weakness
the city names, which establishes that the physical route to F10 is not
speculative. \emph{HydroCanopy}, inspired by frog mucous glands and by the
density gradient of leaves in a canopy, combines solar regulation, ventilation,
thermal inertia and evaporative cooling in one build-up
\parencite{asknature_hydrocanopy}. It is the closest published prior art to our
concept direction, which obliges us to state what we add: the self-regulating
link between load and cooling power (F9), and an effect in the night window.
\todo{For Section~5: p.~128 of the Fachplanung quotes a maximum PET reduction of
\qty{2.3}{\kelvin} in the running text but \qty{1.8}{\kelvin} in Tab.~11 on the
same page. Check the source and say which you use.}

\subsection{System view and project boundary}
\label{sec:systemview}

Figure~\ref{fig:systemview} places the challenge in its system hierarchy. The
hierarchy prevents tunnel vision; the boundary drawn afterwards prevents scope
explosion.

\begin{figure}[htbp]
\centering
\scalebox{0.86}{%
\begin{tikzpicture}
\node[sysbox,minimum width=14.6cm] (super) at (0,4.6)
  {\textbf{SUPER-SYSTEM} --- City / district of Zurich\\
   urban microclimate, street canyon, radiation balance, air flow, urban water cycle};
\node[soibox,minimum width=6.2cm,text width=5.8cm] (soi) at (0,2.4)
  {SYSTEM OF INTEREST\\[1mm]
   \normalfont Heat-regulating outer layer of an\\ existing façade, with the roof\\ as release interface\\
   \normalfont\itshape the thermal interface between\\ building and city};
\node[sysbox,minimum width=3.7cm,text width=3.3cm] (p1) at (-5.5,3.1)
  {\textbf{PARALLEL 1}\\ Stormwater and sewer system};
\node[sysbox,minimum width=3.7cm,text width=3.3cm] (p2) at (5.5,3.1)
  {\textbf{PARALLEL 2}\\ Street and pavement network};
\node[sysbox,minimum width=3.7cm,text width=3.3cm] (p3) at (-5.5,1.4)
  {\textbf{PARALLEL 3}\\ Power grid and building energy system};
\node[sysbox,minimum width=3.7cm,text width=3.3cm] (p4) at (5.5,1.4)
  {\textbf{PARALLEL 4}\\ Residents, transport and public transit};
\node[subbox,text width=2.4cm] (s1) at (-6.1,-0.4) {Optical surface layer};
\node[subbox,text width=2.4cm] (s2) at (-3.05,-0.4) {Heat storage and buffer layer};
\node[subbox,text width=2.4cm] (s3) at (0,-0.4)  {Water capture, store and release};
\node[subbox,text width=2.4cm] (s4) at (3.05,-0.4) {Insulation and substrate};
\node[subbox,text width=2.4cm] (s5) at (6.1,-0.4) {Mounting and fixing};
\node[font=\scriptsize\bfseries] at (0,0.7) {SUBSYSTEMS};
\draw[flowarr] (super) -- (soi);
\foreach \p/\l in {p1/water,p2/radiation,p3/energy,p4/exposure}
  { \draw[{Latex[length=2mm]}-{Latex[length=2mm]},thick,draw=black!45] (\p) -- (soi)
      node[midway,font=\scriptsize,fill=white,inner sep=1pt] {\l}; }
\foreach \s in {s1,s2,s3,s4,s5} { \draw[draw=black!45] (soi) -- (\s); }
\begin{scope}[on background layer]
  \node[draw=zhawblue,dashed,line width=1pt,rounded corners=4pt,
        fit=(soi)(s1)(s5),inner sep=7pt] (bnd) {};
\end{scope}
\node[font=\scriptsize\bfseries,color=zhawblue,anchor=west]
  at ($(bnd.north west)+(0.15,0.30)$) {PROJECT BOUNDARY};
\end{tikzpicture}}
\caption{System view. The dashed line is the project boundary: everything inside
it is designed and analysed by the team. The subsystems are components and
functions, not people --- the trades who install and maintain the system appear
as constraints (Table~\ref{tab:constraints}). What the comparison with each
parallel system contributed is set out in Table~\ref{tab:parallel}.}
\label{fig:systemview}
\end{figure}

\begin{table}[htbp]
\centering
\caption{Project boundary: what is deliberately in and out of the project.}
\label{tab:inout}
\renewcommand{\arraystretch}{1.15}
\small
\begin{tabularx}{\textwidth}{X X}
\toprule
\textbf{IN the project} & \textbf{OUT of the project} \\
\midrule
Passive control of the short- and long-wave radiation balance of the envelope
surface
& Redesign of the building: structure, floor plan, windows, insulation
thickness \\
The self-regulating link between heat load and cooling power
& Any actively controlled system: sensors, valves, pumps, chillers, mechanical
ventilation \\
Heat storage and delayed release within the outermost build-up, incl.\ the night
window
& Seasonal storage, boreholes, heat pumps, thermal networks, electricity
generation \\
Capture, retention and metered release of rainwater within the build-up
& Municipal water supply and sewage; street trees and district-scale greening \\
Mounting on an existing substrate, installability, durability of the function
& Permitting, heritage approval, business model, full life cycle assessment \\
\bottomrule
\end{tabularx}
\end{table}

\paragraph{What the systems view changed.} \emph{The systems view changed our
project scope because} it forced two corrections. First, we carried two systems
of interest --- reducing heat and preventing heat; the hierarchy showed these to
be two \emph{functions} of one system, coupled by the design question, so a
single boundary can be drawn. Second, the comparison with the parallel systems
(Table~\ref{tab:parallel}) revealed the night sky and the discarded rainwater as
the two heat sinks available without energy input, and the pedestrian in the
canyon as the place where the effect must appear. Rainwater is currently drained
away as fast as possible, yet under the \enquote{no potable water} constraint it
is the only water available during a heat wave: seen from our system, a cooling
reserve being discarded. This added F2 and F6--F8, the evening window, and
KPI~5 and~6.

\subsection{Functions and functional diagram}
\label{sec:functions}

Functions are formulated as solution-neutral verbs answering \enquote{what must
the design \emph{do}?}, and stated at the scale of the designed object rather
than the city: the design regulates heat at the surface it occupies, and the
city-scale effect is the \emph{goal}, not the function --- as the function of
polar bear fur is to insulate, not to make the Arctic habitable.

\begin{table}[htbp]
\centering
\caption{Function list. Main function, sub-functions and supporting functions.}
\label{tab:functions}
\renewcommand{\arraystretch}{1.05}
\footnotesize
\begin{tabularx}{\textwidth}{L{0.8cm} L{4.4cm} X}
\toprule
\textbf{ID} & \textbf{Function (verb)} & \textbf{Explanation / taxonomy link} \\
\midrule
\multicolumn{3}{l}{\textit{Main function}}\\
F0 & Regulate heat gain and release at an urban surface & Modify / regulate physical state: manage temperature \\
\midrule
\multicolumn{3}{l}{\textit{Sub-functions --- prevention (day window)}}\\
F1 & Reflect short-wave radiation & Prevent absorption of solar energy in the first place \\
F5 & Shade a surface & Reduce the irradiated area, e.g.\ through surface geometry \\
\midrule
\multicolumn{3}{l}{\textit{Sub-functions --- removal (both windows)}}\\
F2 & Emit long-wave radiation & Release heat to the sky, in particular in the atmospheric window \\
F3 & Transfer heat to ambient air & Convective dissipation at the surface \\
F4 & Store thermal energy temporarily & Buffer and shift the daily peak \\
\midrule
\multicolumn{3}{l}{\textit{Sub-functions --- water path}}\\
F6 & Capture water & Collect rain or atmospheric moisture \\
F7 & Retain water & Store it until it is needed during the heat peak \\
F8 & Release water in a controlled way & Meter evaporation so the reserve lasts the whole heat wave \\
\midrule
\multicolumn{3}{l}{\textit{Sub-function --- the core of the design question}}\\
F9 & Couple response strength to load & Increase removal (F2, F3, F8) as surface temperature rises, without external control \\
\midrule
\multicolumn{3}{l}{\textit{Supporting functions}}\\
F10 & Keep a surface clean & Maintain the optical function against soiling and algae \\
F11 & Resist degradation & Withstand UV, frost, driving rain over the service life \\
F12 & Attach to a substrate & Fix the build-up to an existing wall, installable by ordinary trades \\
\bottomrule
\end{tabularx}
\end{table}

\paragraph{Why F9 is physically plausible.} Every removal path depends on
surface temperature, but not equally strongly. Differentiated with respect to
$T_s$, the self-regulation gradients are, in
\unit{\watt\per\square\metre\per\kelvin}, radiative
$4\varepsilon\sigma T_s^{3} \approx 6$, convective $h_c \approx 10$, and latent
$\propto \mathrm{d}e_s/\mathrm{d}T_s$, which is steep because saturation vapour
pressure roughly doubles per \qty{10}{\kelvin}. The evaporative path is by far
the steepest, which is why transpiration, not radiation, is the mechanism
biology uses when a load must be tracked. \todo{Section~3.2: evaluate the three
derivatives with your own assumptions; this check decides which mechanism
carries F9.}

\begin{figure}[htbp]
\centering
\scalebox{0.86}{%
\begin{tikzpicture}
\node[funcbox,fill=yellow!25,text width=2.1cm] (sun) at (-6.8,2.4)
  {Solar irradiance $G$};
\node[funcbox,fill=zhawblue!12,text width=2.1cm] (rain) at (-6.8,-2.2)
  {Rain / atmospheric moisture};
\node[funcbox,text width=2.4cm] (f1) at (-3.4,3.3) {\textbf{F1/F5} Reflect and shade};
\node[funcbox,fill=red!8,text width=2.4cm] (ts) at (-3.4,1.0)
  {\textbf{F4} Store heat\\[1mm] \itshape state: surface temperature $T_s$};
\node[funcbox,text width=2.4cm] (f2) at (0.5,2.7) {\textbf{F2} Emit long-wave radiation};
\node[funcbox,text width=2.4cm] (f3) at (0.5,1.0) {\textbf{F3} Transfer heat to air};
\node[funcbox,text width=2.4cm] (f8) at (0.5,-0.8) {\textbf{F8} Release water, metered};
\node[funcbox,text width=2.4cm] (f67) at (-3.4,-2.2) {\textbf{F6/F7} Capture and retain water};
\node[funcbox,text width=2.4cm,draw=red!60] (fc) at (0.5,-2.6)
  {\itshape Residual: heat into the building};
\node[funcbox,fill=gray!12,text width=2.0cm] (sky) at (4.5,2.7) {Sky (long-wave sink)};
\node[funcbox,fill=gray!12,text width=2.0cm] (air) at (4.5,0.2) {Street-canyon air};
\node[funcbox,fill=gray!12,text width=2.0cm] (bld) at (4.5,-2.6) {Building interior};
\draw[flowarr] (sun) -- (f1);
\draw[flowarr] (sun) -- (ts) node[midway,below,font=\scriptsize] {absorbed};
\draw[flowarr] (f1.north) .. controls +(1.6,1.2) and +(-1.6,1.2) .. (sky.north west)
  node[midway,above,font=\scriptsize] {reflected};
\draw[flowarr] (ts) -- (f2);
\draw[flowarr] (ts) -- (f3);
\draw[flowarr] (ts) -- (fc);
\draw[flowarr] (f2) -- (sky);
\draw[flowarr] (f3) -- (air);
\draw[flowarr] (f8) -- (air) node[midway,below,font=\scriptsize] {latent};
\draw[flowarr,draw=red!60] (fc) -- (bld);
\draw[wateharr] (rain) -- (f67);
\draw[wateharr] (f67) -- (f8);
% --- the self-regulating loop
\draw[-{Latex[length=2mm]},thick,draw=orange!85!black,dash pattern=on 4pt off 2pt]
  (ts.south east) .. controls +(1.4,-1.4) and +(-1.4,0.6) .. (f8.west);
\node[font=\scriptsize\bfseries,color=orange!85!black,align=center]
  at (-1.5,-1.1) {\textbf{F9}\\ higher $T_s$\\ $\Rightarrow$ stronger removal};
\begin{scope}[on background layer]
  \node[draw=zhawblue,dashed,rounded corners=4pt,fit=(f1)(ts)(f2)(f3)(f8)(f67)(fc),
        inner sep=9pt] (fb) {};
\end{scope}
\node[font=\scriptsize\bfseries,color=zhawblue,anchor=west]
  at ($(fb.north west)+(0.15,0.30)$) {SYSTEM BOUNDARY (heat-regulating outer layer)};
\node[font=\scriptsize,anchor=west] at (-7.0,-3.9)
  {\textcolor{black!65}{\textbf{---}} energy flow \quad
   \textcolor{zhawblue}{\textbf{--\,--}} water flow \quad
   \textcolor{orange!85!black}{\textbf{--\,--}} self-regulating coupling \quad
   \textcolor{red!70}{\textbf{---}} flow to be minimised};
\end{tikzpicture}}
\caption{Solution-neutral functional diagram. Only functions and flows are
shown, no components and no materials. The orange coupling is F9 and is what
distinguishes this project from a fixed high-albedo surface: the state variable
$T_s$ feeds back into the strength of the removal functions. Supporting
functions F10--F12 act on everything inside the boundary and are omitted for
readability.}
\label{fig:funcdiagram}
\end{figure}

\subsection{Biologized questions}
\label{sec:biologized}

The functions are translated into questions to nature, which form the search
strategy for Section~3.1. The challenge brief asks one question directly --- what
can be learned from desert plants and animals about passive cooling --- and the
answer shapes where we looked: the organisms most challenged by our problem, and
unfazed by it, live where solar load is extreme and water is scarce. That is why
the models of Section~\ref{sec:discovering} are disproportionately desert
organisms, and why the two that transfer best are those whose strategy is a
structure rather than a metabolic process.

\begin{itemize}
  \item How does nature use the heat it receives to get rid of that same heat?
  \item How does nature make a response stronger when the load is stronger,
        without a controller?
  \item How does nature release excess heat to the night sky?
  \item How does nature reflect intense solar radiation without pigments that
        bleach, and without a mirror finish?
  \item How does nature capture, store and meter out scarce water so that a
        reserve lasts a whole dry period?
  \item How does nature buffer a daily temperature extreme instead of resisting
        it?
  \item How does nature keep a functional surface clean without cleaning it?
\end{itemize}

\subsection{Performance criteria and KPIs}
\label{sec:kpi}

The KPIs are derived from the functions, interfaces and constraints identified
above, not invented independently of them. They are used twice: to guide the
engineering checks in Section~3.2 and to compare the concepts in Section~3.3.

\begin{table}[htbp]
\centering
\caption{Performance criteria. Function $\rightarrow$ KPI $\rightarrow$ unit
$\rightarrow$ target. Reference case: an existing Zurich façade with a solar
reflectance of about $0.3$, the value the city uses for un-greened façades
\parencite[114]{stadtzuerich2020}.}
\label{tab:kpi}
\renewcommand{\arraystretch}{1.05}
\footnotesize
\begin{tabularx}{\textwidth}{L{0.4cm} L{4.3cm} L{1.4cm} X}
\toprule
\textbf{\#} & \textbf{Function $\rightarrow$ KPI} & \textbf{Unit} & \textbf{Target / benchmark} \\
\midrule
1 & F9 --- self-regulation gradient
    $\partial P_\mathrm{cool}/\partial T_s$
  & \unit{\watt\per\square\metre\per\kelvin}
  & Clearly positive and larger than for a fixed high-albedo surface.
    \textbf{The defining criterion of this project.} \todo{set a target after
    the derivative check} \\
2 & F1 --- solar reflectance $\alpha_\mathrm{sol}$ (300--2500\,nm)
  & --
  & $\geq 0.80$, the city's threshold for \enquote{high albedo}
    \parencite[128]{stadtzuerich2020}; baseline $0.3$, ultra-white laboratory
    paints $\approx 0.98$ \parencite{li2021}. \\
3 & F0 --- peak surface temperature reduction $\Delta T_s$, clear sky, 14:00
  & K
  & $\geq \qty{15}{\kelvin}$. \todo{confirm before claiming it} \\
4 & F0 --- PET reduction at \qty{2}{\metre} above ground, 14:00
  & K
  & Beat HA~11's \qty{-1.0}{\kelvin} median for a bright \emph{façade}; the
    \qty{-1.2}{\kelvin} for a roof is context, not our benchmark
    \parencite[Tab.~11]{stadtzuerich2020}. The metric the client uses. \\
5 & F2, F4 --- net cooling power at 02:00, clear sky; heat discharged
    18:00--06:00
  & \unit{\watt\per\square\metre}; \unit{\watt\hour\per\square\metre}
  & $\geq \qty{40}{\watt\per\square\metre}$ instantaneous
    \parencite{raman2014,zhai2017}, and a night-time air temperature effect
    better than HA~11's \qty{-0.1}{\kelvin}. \\
6 & F6--F8 --- potable water use; days of autonomy on harvested rainwater
  & \unit{\litre\per\square\metre\per\day}; d
  & Potable water $= 0$ (hard constraint); autonomy $\geq$ \qty{5}{days}, so the
    reserve survives the design heat wave. \todo{derive the storage from
    \qty{2.45}{\mega\joule\per\kilo\gram}} \\
7 & F10/F11 --- retention of the cooling function after soiling and weathering
  & \unit{\percent} of initial
  & $\geq \qty{90}{\percent}$ after three years. Soiling is the weakness the
    city names \parencite[128]{stadtzuerich2020}, and aged cool roofs commonly
    lose a significant fraction of their reflectance
    \parencite{santamouris2014}. \\
\bottomrule
\end{tabularx}
\end{table}

\paragraph{Two conditions on how these are read.} First, the \emph{reset
condition}: KPI~1 and~5 are read per day, but over the five-day case the surface
temperature and the state of charge of the store at 06:00 must not drift upwards
from one day to the next. A concept that cools well on day one and saturates by
day three fails the design case even if every single-day KPI is met.
\todo{report the five-day trend, not only day one} Second, KPI~3 measures a
\emph{surface} temperature and KPI~4 an effect on \emph{people}: a
high-reflectance surface can be \qty{20}{\kelvin} cooler than a dark one while
the effect on district air temperature is smaller by roughly two orders of
magnitude \parencite{santamouris2014,akbari2009}. Conflating them is the most
common error in urban-heat-island projects, so they are reported separately
throughout. The targets are aspirational benchmarks set during scoping ---
hypotheses, not results --- and each is confirmed, revised or declared unverified
in Section~3.2. \todo{verify every cited benchmark against its primary source}

\paragraph{Three changes to the team's original KPI list.} Installed cost and
embodied greenhouse gas were an eighth KPI; they are now constraints
(Table~\ref{tab:constraints}), because a KPI must follow from a \emph{function}
and neither does --- which also brings the list to the seven the task asks for.
\emph{Urban heat island intensity} was dropped as a selection KPI, being a
district-scale quantity a single retrofit cannot move measurably, and is
replaced operationally by KPI~4. \emph{Sustainability cost} was not measurable
as stated and became the embodied-emissions constraint, which has a unit, a
benchmark and a decision rule.

\subsection{Relevant Life's Principles}
\label{sec:lifesprinciples}

The integrate step asks which principles are \emph{relevant}, not how well the
design achieves them; the assessment belongs to
Section~\ref{sec:sustainability} and is made there against
Table~\ref{tab:lpassess}. Each is named here with the sub-principle claimed.
\textbf{Be locally attuned and responsive} --- \emph{signal, antennae and
response match} --- is the principle the design question is built on, and F9 is
it expressed as a function. \textbf{Be resource efficient} contributes
\emph{use low energy processes}, which is the hard energy constraint, and
\emph{one design meets more than one need}, a tie-breaker in
Section~\ref{sec:evaluation}. \textbf{Adapt to changing conditions} contributes
\emph{adapt to temporal change} (a five-day heat wave, not one afternoon),
\emph{withstand disturbance} (frost, driving rain, UV) and \emph{maintain
integrity through self-renewal} (KPI~7). \textbf{Use life-friendly chemistry}
contributes \emph{a safe subset} and \emph{harmless breakdown products}, which
is not decorative here: whatever leaches from the surface enters rainwater the
concept evaporates at pedestrian level. \textbf{Integrate development with
growth} is claimed only as \emph{built to shape} --- \emph{creates more
opportunities for life} is deliberately \textbf{not} claimed, because a
permanently damp textured surface invites the algal growth KPI~7 penalises. And
\textbf{evolve to survive} is \textbf{not claimed for the artefact} at all: a
façade layer does not evolve across generations, so it is claimed for the
\emph{process} only.

The umbrella principle, \emph{create conditions conducive to life}, requires a
design to be \emph{optimised rather than maximised}. That is a decision here,
not a platitude: ultra-white coatings reach about $0.98$ \parencite{li2021}, but
KPI~2 targets the city's own $\geq0.80$, because further reflectance costs
glare, winter gains and photovoltaic yield.

\paragraph{Quality check.} System of interest clearly defined? Yes, and
singular. Project boundary manageable? Yes, Table~\ref{tab:inout}; seasonal
storage and electricity generation are explicitly out. Functions
solution-neutral? Yes, verbs only, with no components in
Table~\ref{tab:functions} or Figure~\ref{fig:funcdiagram}. Did the systems view
reveal additional functions or constraints? Yes: F2 and F6--F8 from the sky and
the discarded rainwater, the evening window from the residents, and the
installation and maintenance constraints. Most important functions measurable?
Yes, KPIs~1--7, each with a unit. Is the \enquote{How might we\ldots?} question
open to different solutions? Yes: evaporative, radiative and geometric
strategies all remain possible, and Section~3.2 must choose between them on
evidence.
